\section{Future Work}

DIAL currently implements the core path: nodes advertise capabilities, discover
peers, route queries, and return answers through the selected node. Future work
can improve this prototype in two ways. Some directions address limits observed
during evaluation, such as candidate visibility, timeout-heavy failures, and
policy comparison. Others extend the project toward a richer decentralized AI
system, with several nodes working together, better trust signals, privacy
constraints, and incentives for open participation. The ideas below are
therefore both practical improvements to the current prototype and possible
extensions of the broader DIAL vision.

\subsection{Multi-Node Query Execution}

The current router selects one node. A natural extension is multi-node
execution, where a request can involve several specialized peers. One version
is top-\(k\) execution: send the same query to several high-ranking nodes, then
aggregate the answers. Another version is staged execution: route a request
through a small sequence of nodes, for example planning first, then coding,
then writing or verification. This could improve robustness and quality in
several ways:

\begin{itemize}
  \item compare answers and select the most consistent one;
  \item ask answer evaluator nodes to rate responses;
  \item combine specialized answers, for example one math model and one writing
  model;
  \item retry automatically when one selected peer times out;
  \item expose answer diversity to the user.
\end{itemize}

Multi-node execution fits the decentralized direction. A network of smaller
specialized models may be more useful than one local model when their strengths
are complementary. The trade-off is coordination cost: more nodes can improve
quality or robustness, but they also add latency and failure points. This
connects to LLM-Blender and Mixture-of-Agents, which show that ranking, fusing,
or layering multiple model outputs can improve quality
\citep{jiang2023llmblender,wang2024mixtureofagents}.

\subsection{Aggregation and Answer Evaluation}

Once several nodes answer, the system needs aggregation. Aggregator nodes could
be regular DIAL nodes whose role is to fuse several responses into one final
answer. They could select the strongest parts of each answer, remove
contradictions, summarize long responses, or produce a final response with
citations and confidence notes.

The network may also need answer evaluator nodes. Their role would be different
from capability evaluators: instead of estimating whether a node is generally
good at math or writing, they would judge whether a specific answer was useful,
correct, complete, or well-grounded. This could happen during multi-node
execution, or periodically as a quality check to see whether selected nodes are
still producing good answers. Verifier-based systems such as LLM-as-a-Verifier
are relevant because they turn answer checking into a reusable component
\citep{kwok2026llmverifier}.

This creates a new routing problem: the best set of nodes is not always the top
\(k\) individually. The router may want diversity across capabilities, model
families, latency classes, trust domains, and node roles. Matching then becomes
a problem of choosing a good group of answer nodes, evaluator nodes, and
possibly an aggregator node.

\subsection{Capability Profile Calibration}

The current prototype assigns capability profiles from configured or simulated
scores. A stronger system should learn these profiles from evidence. Each node
could periodically answer a small evaluation suite covering math, programming,
writing, summarization, research, planning, creative tasks, and general
questions. Existing work on user-need taxonomies and language-model evaluation
could help design such a suite without reducing capability to one global score
\citep{shelby2025tuna,liang2022helm}. The results could be graded by tests,
human feedback, or verifier models, then used to update the node profile over
time.

This could be handled by capability evaluator nodes: regular DIAL nodes whose
task is to test other nodes and publish signed or local evaluation results. It
would make advertised capabilities less dependent on manual role hints or
self-declared profiles. It would also give reputation a better input: routing
could learn not only whether a node is reachable, but whether it actually
performs well on the tasks it claims to support.

\subsection{Reputation}

Reputation could measure whether a node is useful over time. Signals include
uptime, latency, successful completion, user feedback, verifier scores, and
consistency with other nodes. A first implementation can maintain local scores;
later versions can publish signed reputation events or use verifiable logs.

Reputation is also a defense mechanism. In an open network, nodes may advertise
capabilities they do not have. Routing should learn from observed behavior
rather than trust advertisements blindly. Verifiable inference mechanisms such
as VeriLLM point toward stronger versions, where selected execution claims can
be checked publicly
\citep{wang2026verillm}.

\subsection{Incentives and Pricing}

If the network becomes a market, providers need incentives. A future
\texttt{NodeProfile} could include price metadata such as
\texttt{price\_per\_query}, accepted payment methods, free quotas, or priority
fees. Routing could then include cost in the matching objective, instead of
choosing only by capability and operational health.

Blockchain or payment-channel mechanisms become relevant here. The discovery
plane can remain peer-to-peer, while the payment layer records payments,
escrows, staking, or rewards. This keeps the architecture close to BitTorrent's
lesson: incentives help open networks stay robust, but should serve the
resource exchange rather than replace discovery \citep{cohen2003bittorrent}.
It also fits the broader PPAI direction, where matching requests to providers
can include more than technical capability, such as cost and participation
rules \citep{wang2026ppai}.

\subsection{Faster Failure Detection}

The dropout and high-latency benchmarks show that failure handling still needs
work. The prototype often still returns an answer when many peers disappear,
which is a good sign, but weaker candidate visibility and slow failed paths
reduce quality and increase tail latency. Future versions could add:

\begin{itemize}
  \item shorter adaptive timeouts based on recent latency;
  \item profile expiration and explicit unavailability propagation;
  \item background health checks;
  \item multiple candidate reservations before forwarding;
  \item fallback to local execution when remote confidence is low.
\end{itemize}

These changes would likely improve quality and p95 latency when peers are
slow, stale, or unavailable. They also align with agent-discovery work that
treats service readiness and changing peer availability as separate but
interacting discovery problems
\citep{dazzi2026usableagentdiscovery,guo2026agentdiscovery}.

\subsection{Load-Aware Request Spreading}

The current policy framework can compare capability quality, freshness,
failures, and latency, but it does not yet include an explicit load-balancing
term. In a larger network, this matters because the best node on paper should
not necessarily receive every request. If too many users choose the same strong
peer, that peer may become slower or unavailable, while other good-enough peers
remain unused.

A future router could therefore track simple load signals, such as recent
selections, active requests, queue length, or provider-declared capacity. The
router would still prefer capable nodes, but it could spread requests across
several suitable peers instead of always selecting the highest score. This fits
the PPAI idea of matching requests to providers under several constraints:
quality matters, but so do availability, cost, trust, and network health
\citep{wang2026ppai}. Existing routing systems such as RouteLLM and the vLLM
Semantic Router also support this view of routing as a policy decision over
multiple signals, even if they do not target the same peer-to-peer setting
\citep{ong2025routellm,vllmsemanticrouter2026}.

\subsection{Discovery Coverage and Adaptive Recommendations}

The evaluation suggests that several weaknesses come from limited candidate
visibility. A first improvement direction is therefore to tune the discovery
layer itself: DHT provider lookup breadth, provider refresh timing, profile
refresh behavior, bootstrap warmup, and how many providers are kept per
capability. These are engineering parameters rather than new research ideas,
but they matter because the router can only score candidates it has discovered.

The current recommendation mechanism is also conservative. The entry node
refreshes its candidate pool through the DHT, asks a few eligible peers for
recommendations when needed, and stops. This keeps routing predictable, but it
also makes discovery shallow. In a larger or uneven network, useful nodes may
be beyond the peers initially found by the DHT or known by the entry node.
This matters especially for mixed-capability queries, where the scoring formula
needs enough candidates across several capabilities before it can make a good
trade-off.

A future router could use adaptive multi-hop recommendation expansion. When
candidate quality is low, the score gap is small, or the entry node has too
few good options, it could ask recommended peers for their own recommendations.
The goal would be to discover further into the network without turning routing
into an uncontrolled search. The expansion should remain bounded:

\begin{itemize}
  \item set a maximum hop depth, for example two or three recommendation hops;
  \item set a total recommendation budget per query;
  \item stop early when a reachable candidate passes the routing threshold;
  \item stop when new hops return mostly duplicates or already failed peers;
  \item penalize longer-path candidates unless their profile is clearly better;
  \item cache useful paths briefly, but expire them aggressively when peers
  change or disappear.
\end{itemize}

This trades extra discovery cost for a wider view of the network when one-hop
routing is not enough. It is a middle ground between shallow DHT lookup and
uncontrolled flooding: discovery can expand when useful, but still remains
bounded.

\subsection{Privacy-Aware Routing and Trust Domains}

Not all queries should be routed to any peer. Some prompts contain private,
institutional, or regulated information. The matching model can support trust
constraints: route only to local peers, known identities, institutional peers,
or peers with a privacy policy. This makes privacy a routing decision: some
tasks may stay inside a trusted boundary, while less sensitive tasks may use a
broader public network.

Privacy is harder in a public peer-to-peer network. DIAL should not assume that
a raw prompt remains private once it is sent to an unknown execution node. The
current design limits exposure during discovery: DHT lookup and recommendations
exchange capabilities, profiles, and peer IDs, while the prompt is sent only to
the selected node for execution. Still, the entry node and selected node can see
the prompt.

This also opens the possibility of private or hybrid DIAL networks. An
individual, lab, or organization could run a small private network of local
specialized nodes instead of relying on one large centralized model. Such a
network could keep sensitive prompts inside a trusted boundary while still
benefiting from specialization and aggregation. A hybrid version could route
private tasks locally and use a broader public network only when policy allows
it. This may improve privacy, specialization, and robustness, but it would need
careful policies to avoid unnecessary latency or leakage. This is consistent
with distributed inference systems such as Petals, where shared compute is
useful but privacy remains an important deployment concern
\citep{borzunov2023petals}.

A future version could make this explicit with privacy modes such as
\texttt{local-only}, \texttt{trusted-only}, \texttt{institutional}, or
\texttt{public}. It could also reduce exposure by classifying locally before
discovery, routing on capability metadata, redacting sensitive fields when
possible, or sending only non-sensitive sub-tasks to public peers. More advanced
approaches, such as privacy-preserving routing or privacy-preserving
decentralized serving, are explored by PPRoute and PlanetServe
\citep{wu2026privacyrouting,fang2026planetserve}. These approaches show that
privacy can be a design direction for decentralized LLM routing, but it is not
automatic and usually adds protocol or infrastructure cost.

\subsection{Security and Identity}

A future open network would need stronger identity guarantees than the current
local prototype. Today, a node can publish a profile saying which capabilities
it offers, but the prototype does not verify that this profile was signed by a
stable identity or that the same actor is not creating many disposable nodes.
This matters because reputation, payment, and capability evaluation only work
if the network can attach past behavior to a persistent node identity.

A stronger design could require signed profiles, timestamps or freshness
checks, and a clearer binding between a peer ID, public key, and advertised
profile. It could also make fake or duplicated node identities harder to use,
for example by linking reputation to longer-lived identities or requiring some
form of verification before a node is trusted for sensitive or paid work. The
lightweight discovery layer should therefore be paired with stronger
authentication beyond local experiments. Lattica's cross-NAT focus is relevant
because an open network cannot assume that all peers are directly reachable
\citep{yang2025lattica}.

\subsection{Richer Matching Algorithms}

The current vector model is deliberately simple. A later router could replace
the current score with a broader utility objective:

\[
  U(n,q) =
  w_q Q(n,q) + w_r Rep(n) - w_p Price(n) - w_l Latency(n) - w_\rho Risk(n).
\]

This would let the system choose, for example, a cheaper node above a quality
threshold, a high-reputation node for sensitive requests, or a diverse group of
nodes when aggregation is useful. Future matching can also use:

\begin{itemize}
  \item semantic embeddings of queries and model descriptions;
  \item learned routing from historical outcomes;
  \item load-aware and energy-aware scheduling;
  \item price-quality trade-off curves;
  \item task decomposition across multiple specialist nodes;
  \item user-specific preferences for speed, quality, privacy, or cost.
\end{itemize}

This is where the project connects most strongly to PPAI. The network does not
need to hard-code one matching algorithm; it can expose data and interfaces
for competing policies. PPAI, UniRoute, RouteProfile, and vLLM Semantic Router
suggest richer ways to represent tasks, models, policies, and constraints
\citep{wang2026ppai,jitkrittum2025uniroute,xu2026routeprofile,vllmsemanticrouter2026}.

\subsection{Beyond the Prototype}

If DIAL were extended beyond the current prototype, the natural order would be
to make the network more reliable, add trust and reputation so routing can
learn from experience, and then make participation more attractive for people
or organizations that run useful nodes. Incentives for providers could mean
visibility for reliable nodes, priority rules, credits, payments, or other
rewards for sharing compute and specialized models. A later system would look
less like a single application and more like an open protocol for AI service
exchange.
